\paragraph{E.3m (Proved: the exact integral Walker reformulation of the
general family converse).}
Fix a smooth projective complex variety $X$ and a codimension $i$.
The following assertions are equivalent:
\begin{enumerate}
\item For every abelian variety $A$ and every normalized
      $Z\in\mathrm{CH}^i(A\times X)$, the vanishing
      $a({}^tZ)=0$ implies $u_Z=0$ integrally.
\item The natural surjection
      $P^i(X)\to J_W^i(X)(\mathbb C)$ is an isomorphism.
\item $N^i(X)=\ker\psi_W$ in $\mathrm{CH}^i(X)_{\mathrm{alg}}$.
\end{enumerate}
This is a reduction, not an assertion that these equivalent
conditions hold for every $X,i$. In particular the remaining
general gate is the exact Walker-kernel equality, rather than
the false integral ordinary Abel--Jacobi equality in Lead E4.

We write $J=J_W^i(X)$, $P=P^i(X)$, and
$\overline\psi:P\twoheadrightarrow J(\mathbb C)$.
The equivalence of (2) and (3) follows from surjectivity of
$\psi_W$ and E.3g's inclusion $N^i\subset\ker\psi_W$.
Condition (2) implies (1) by E.3i. We prove the other direction.

\emph{Choose one family covering the Walker Jacobian.}
Assume (1). Every algebraically trivial cycle is an integral
point of an abelian family by E.3g. The corresponding regular
Walker maps have abelian-subvariety images which generate $J$.
A finite sum of these images is all of $J$: otherwise one can
successively enlarge their dimensions, which is impossible
after $\dim J$ strict enlargements. Thus a product of finitely
many such parameter varieties gives an abelian family
\[
 u_B:B(\mathbb C)\longrightarrow P,\qquad
 \phi_B=\overline\psi\,u_B:B\twoheadrightarrow J.
\]
If $J=0$, (1) immediately kills every family and hence $P$;
we may therefore assume $\dim J>0$.

Let $B_0=(\ker\phi_B)^0$. The restricted family on $B_0$
has zero Walker map and therefore zero cohomological action
by E.3b--i. Condition (1) gives $u_B(B_0(\mathbb C))=0$.
Consequently $u_B$ factors through $D=B/B_0$, and
$D\to J$ is an isogeny. Its abstract kernel on $P$ is a
subgroup of the finite group $\ker(D\to J)(\mathbb C)$.
Over $\mathbb C$ this subgroup is the point group of a finite
subgroup scheme $H\subset D$. Put $E=D/H$. We obtain
\[
 u_E:E(\mathbb C)\hookrightarrow P,
 \qquad \lambda:E\twoheadrightarrow J,
 \qquad \overline\psi\,u_E=\lambda,
\]
where $\lambda$ is an isogeny and $u_E$ is injective. At this
point $u_E$ is a homomorphism of groups; we have not assumed
that a universal integral cycle exists on $E\times X$.

\emph{Every abelian-family image is already this one.}
For any other abelian family $u_A:A(\mathbb C)\to P$, take
the sum family on $B\times A$ and repeat the preceding
construction. Its Walker map is still surjective. We get an
abelian variety $E'$ isogenous to $J$ and an injection
$u_{E'}:E'(\mathbb C)\hookrightarrow P$ whose image is
$u_B(B(\mathbb C))+u_A(A(\mathbb C))$.
The inclusion $B\to B\times A$ induces an algebraic
homomorphism $B\to E'$. It kills the kernel of $B\to E$,
because its composition with the injective map $u_{E'}$
is $u_B$. Quotient descent therefore gives an algebraic
homomorphism $E\to E'$. It is injective on complex points,
again by the injections into $P$. Both abelian varieties have
dimension $\dim J$, so this homomorphism is an isomorphism.
It follows that $u_A(A(\mathbb C))\subset u_E(E(\mathbb C))$.
Using E.3g once more gives
\[
                         P=u_E(E(\mathbb C)).
\]
Moreover, the preceding quotient construction shows that
every composite $u_E^{-1}u_A:A\to E$ is an algebraic
homomorphism, not merely an abstract homomorphism on points.

The kernel of $\overline\psi$ is now the finite group
$\ker\lambda(\mathbb C)$. But it is also
$\ker\psi_W/N^i(X)$, a divisible group by
\citet[Theorem 2.3, author p.~12]{cyACMVWalker2022}.
A finite divisible group is zero. Hence $\lambda$ is an
isomorphism and (2) follows. This already completes the
equivalence; no universality assumption on the Walker map
has entered the argument.

\emph{Regularity of the resulting quotient, also checked directly.}
Let
\[
 q:\mathrm{CH}^i(X)_{\mathrm{alg}}longrightarrow P
                       \xrightarrow{u_E^{-1}}E(\mathbb C).
\]
For a pointed smooth projective parameter variety $(T,t_0)$
and a correspondence $W$, its point-to-cycle map modulo $N$
factors integrally through $\operatorname{Alb}(T)(\mathbb C)$,
because $W_*$ kills the Albanese kernel by the definition of
$N$. Denote the resulting homomorphism to $E(\mathbb C)$ by
$v$. E.3a constructs an integer $m>0$ and an abelian family
on $\operatorname{Alb}(T)$ whose map to $P$ is $m$ times
the original one. Thus $mv$ is an algebraic homomorphism
$f:\operatorname{Alb}(T)\to E$.

We check that $v$ itself is algebraic. For every
$a\in\operatorname{Alb}(T)[m](\mathbb C)$,
$f(a)=mv(a)=v(ma)=0$. In characteristic zero this means
that $f$ kills the whole finite group scheme
$\operatorname{Alb}(T)[m]$. Hence $f$ factors algebraically
through $[m]$ as $f=g\circ[m]$ for a homomorphism
$g:\operatorname{Alb}(T)\to E$. The difference $v-g$ takes
values in $E[m](\mathbb C)$; its source is divisible, so this
finite-valued homomorphism is zero. Therefore $v=g$.
Composing with the pointed Albanese map proves regularity
for projective parameter varieties.

For a smooth quasi-projective parameter variety, choose a
smooth projective compactification and close the correspondence;
the refined fibres on the original open are unchanged. This
reduces regularity to the preceding case. General smooth
finite-type parameter schemes can be covered by such opens,
and the resulting morphisms agree because their complex-point
maps agree and the target is separated. Thus $q$ is a regular
homomorphism in the usual sense.

We have, in particular,
\[
                    \psi_W=\lambda\circ q.
\]
The direct regularity check is consistent with the alternative
isogeny-initiality proof from the same Theorem 2.3, but that
alternative is unnecessary after the finite-divisible-kernel
argument. The assertion that the original family converse
holds for every $X,i$ remains unproved here.
